\chapter{系统模型与问题定义}
\section{时隙、任务与调度语义}
系统包含 $K$ 架固定高度 UAV、$S$ 个可全频复用的子载波和 $N$ 个地面用户。时隙 $t$ 的到达集合记为 $\mathcal A_t$。任务 $i$ 具有唯一标识，并由
$q_i=(D_i,C_i,T_i,f_i^{\mathrm{loc}})$ 表示。调度项
$e_i=(i,u_i,s_i,r_i)$ 显式记录任务、用户、实际服务 UAV、子载波和 NOMA 角色。任务至多出现一次，每个 $(u,s)$ 块最多包含两项。

\section{信道、全频复用与 SIC}
令 $h_{i,k}$ 为用户 $i$ 到 UAV $k$ 的信道功率增益，$p_i$ 为发射功率。同一子载波上来自其他 UAV 小区用户的干扰为
\begin{equation}
I_{k,s}^{\mathrm{cross}}=\sum_{k'\ne k}\sum_{j\in\mathcal G_{k',s}}p_jh_{j,k}.
\end{equation}
对 UAV $k$ 上的强用户 $a$ 和弱用户 $b$，接收功率满足
$p_ah_{a,k}\ge p_bh_{b,k}$，其 SINR 为
\begin{align}
\gamma_a&=\frac{p_ah_{a,k}}{p_bh_{b,k}+I_{k,s}^{\mathrm{cross}}+\sigma_s^2},\\
\gamma_b&=\frac{p_bh_{b,k}}{\epsilon p_ah_{a,k}+I_{k,s}^{\mathrm{cross}}+\sigma_s^2}.
\end{align}
只有 $\gamma_a\ge\Gamma_{\mathrm{SIC}}$ 时配对才可行，速率为
$R_i=B_s\log_2(1+\gamma_i)$。候选跨 UAV 后必须使用新的实际接收机增益重新排序角色，不能沿用原小区角色。

\section{任务时延和 CPU 约束}
卸载比例为 $\alpha_i$，MEC CPU 为 $f_i$，结果数据比为 $\rho$。本地和边缘分支分别为
\begin{align}
L_i^{\mathrm{loc}}&=\frac{(1-\alpha_i)C_i}{f_i^{\mathrm{loc}}},\\
L_i^{\mathrm{edge}}&=\frac{\alpha_i(1+\rho)D_i}{R_i}+\frac{\alpha_iC_i}{f_i}.
\end{align}
并行分支的完成时延为 $L_i=\max(L_i^{\mathrm{loc}},L_i^{\mathrm{edge}})$。每架 UAV 满足 $\sum_{i:u_i=k}f_i\le F_k$。

\section{优化目标与约束}
主目标是最大化满足 $L_i\le T_i$ 的任务数。为使搜索确定且不牺牲主目标，使用词典序目标
\begin{equation}
\Phi=(\text{完成数},\ \text{总裕量},\ -\text{跨 UAV 干扰}).
\end{equation}
所有结构、角色、SIC 和 CPU 约束均为硬约束。完整枚举仅在不超过五个用户的小实例上作为经验最优参考，不构成一般规模全局最优保证。
